\subsection{Architecture}

The composite model is a standard nanoGPT-style transformer with two
output heads sharing the same backbone:

\begin{itemize}
  \item \texttt{head\_ar}: standard LM head, applied to causal-masked
        hidden states, predicting the next token.
  \item \texttt{head\_diff}: mask-prediction head, applied to
        bidirectionally-attended hidden states, predicting the
        identity of \texttt{<MASK>}-position tokens
        (MDLM-style~\citep{sahoo2024simple}).
\end{itemize}

A \texttt{mode} flag at forward time selects causal vs bidirectional
attention and the corresponding head. Vocabulary is GPT-2 BPE
\citep{radford2019gpt2} extended with a single \texttt{<MASK>} token
(total $V=50{,}258$).

We instantiate the architecture at five scales:

\begin{table}[h]
\centering
\begin{tabular}{lrrrr}
\toprule
Scale & $d_{\text{model}}$ & $L$ & $H$ & Params \\
\midrule
60M  & 512  & 8  & 8  & $\sim 60$M  \\
120M & 640  & 10 & 10 & $\sim 120$M \\
200M & 896  & 12 & 16 & $\sim 200$M \\
250M & 960  & 12 & 16 & $\sim 250$M \\
300M & 1024 & 14 & 16 & $\sim 300$M \\
\bottomrule
\end{tabular}
\caption{Composite model sizes.}
\end{table}

\subsection{Training Recipe}

The joint loss is $\mathcal{L} = \alpha\,\mathcal{L}_{\text{AR}} +
(1-\alpha)\,\mathcal{L}_{\text{diff}}$, with $\alpha$ scheduled linearly
from $1.0$ at training start (pure-AR warm-up) to $0.5$ at training
end. At each step we sample a window from GSM8K-train and apply
either the AR loss (token cross-entropy under causal masking) or the
diffusion loss (cross-entropy on masked positions, with mask ratio
drawn $\text{Uniform}(0.1, 0.9)$), with the AR/diff choice weighted by
$\alpha$.

Batch size $B = 8$; sequence length $T = 256$; optimizer AdamW
($\beta_1=0.9$, $\beta_2=0.95$, weight decay $0.1$); peak LR
$6\times 10^{-4}$ with cosine decay to $10\%$, $200$-step linear
warm-up; bf16 mixed precision. Training duration is $3{,}000$ steps by
default ($\sim 3$ epochs of GSM8K-train); we also run $6{,}000$- and
$10{,}000$-step variants where noted.

\subsection{Baselines}

\textbf{B1 (AR-only)}: identical architecture, $\alpha \equiv 1.0$,
only \texttt{head\_ar} receives gradient.

\textbf{B2 (Diffusion-only)}: identical architecture, $\alpha \equiv
0.0$, only \texttt{head\_diff} receives gradient.

\textbf{B3 (Paired-separate)}: two smaller transformers
($d_{\text{model}}=384$, $L=6$) with the same total parameter count as
the composite, trained \emph{separately} on AR loss and diffusion loss
respectively. Inference is two-pass: one model generates the AR
prefix, the other diffuses the suffix --- the toy-scale analogue of
the inference-pipeline approach in our earlier work on frozen
LLaDA+Qwen substrates.

\subsection{Evaluation Protocol}

\textbf{AR-axis NLL} (\arNLL): for each held-out GSM8K problem, we
compute the per-token cross-entropy on the answer region given the
prompt under causal-mode forward. Held-out chunk: $100$ problems
starting at training-set index $7{,}000$.

\textbf{Diffusion-axis NLL} (\diffNLL): on the same held-out chunk we
mask $50\%$ of answer-region tokens uniformly at random and compute
NLL on masked positions under bidirectional forward. This is the
\emph{mask-fill} task pure-diffusion was trained for; we use it as
the diffusion-axis yardstick.

\textbf{Mode-switching inference}: for downstream accuracy on GSM8K-dev
($N=50$), we evaluate four inference modes per checkpoint:
\texttt{ar\_only} (greedy AR decode); \texttt{mode\_switch\_96\_32}
(AR for 96 tokens, re-mask last 32, $16$-step diffusion-revise);
\texttt{mode\_switch\_64\_32} (analogous); \texttt{paired\_64\_64} (AR
$64$, diffusion-fill next $64$).

\textbf{OOD prose} (D4): AR-mode NLL on WikiText-2-raw,
The Pile-arxiv, and an OpenWebText held-out chunk.

\textbf{Calibration} (D5): Expected Calibration Error (ECE) on AR-mode
predictions, $10$-bin reliability diagram.

\textbf{Compute-matched control}: pure-diffusion baseline trained for
$2\times$ the composite's training steps ($6{,}000$ steps when composite
is at $3{,}000$). This rules out the "more gradient signal" explanation
for any composite diffusion-axis win.

\subsection{Substrate Choice}

We use GSM8K-train ($7{,}473$ problems, $\sim 4$M tokens). This is a
deliberately small dense reasoning substrate: at toy scale, AR
accuracy on GSM8K-dev is $\le 6\%$, so accuracy is binomial-noise
dominated. We rely on continuous NLL metrics throughout and treat
accuracy as a secondary downstream check.
